%==============================================================================
% CHAPTER 13: Deciphering Scrolls with Tomography
%==============================================================================
\chapter{Deciphering Scrolls with Tomography: A Training Experiment}
\label{ch:tomography}
\cite{foschiatti2025tomography}

\begin{flushright}
\textit{Sonia Foschiatti, Axel Kittenberger, Otmar Scherzer}\\
Communicated by \textit{Notices} Associate Editor Krystal Taylor
\end{flushright}

\epigraph{``Virtual unwrapping uses software to reveal text that was hidden for centuries, thereby making it accessible to mankind.''}{--- Foschiatti, Kittenberger, Scherzer}

\begin{abstract}
X-ray computed tomography (CT) enables non-destructive reading of damaged scrolls. This chapter presents a complete didactic pipeline: from simulated X-ray projections (using visible light and transparent films) through filtered back-projection, \index{filtered back-projection}manual segmentation, and Maximum Intensity Projection (MIP) to readable text. Includes Arduino-controlled rotation stage, Python/MATLAB reconstruction code, and the Archeolab software framework, \cite{foschiatti2025tomography}.
\end{abstract}

%==============================================================================
\section{Introduction: The Virtual Unwrapping Problem}
\label{sec:tom:intro}

Damaged scrolls (Herculaneum \index{Herculaneum scrolls}papyri, En-Gedi scroll, medieval manuscripts) cannot be physically opened. X-ray CT provides a solution: \cite{foschiatti2025tomography}
\begin{enumerate}
    \item \textbf{Scan}: Rotate scroll, collect projections (radiographs)
    \item \textbf{Reconstruct}: Filtered back-projection $\to$ 3D volume
    \item \textbf{Segment}: Identify individual layers in the volume
    \item \textbf{Texture/Map}: Assign intensity values to surface points
    \item \textbf{Flatten}: Map 3D surface $\to$ 2D image
    \item \textbf{Merge}: Combine segments into complete text
\end{enumerate}

This chapter implements the full pipeline as an \textbf{educational laboratory} using visible light instead of X-rays.

%==============================================================================
\section{The Mathematical Model: Radon Transform} \cite{foschiatti2025tomography}
\label{sec:tom:radon}

Let $f(x,y)$ be the attenuation coefficient in a slice. A parallel-beam projection at angle $\theta$:
\[
p_\theta(s) = \int_{-\infty}^{\infty} f(s\cos\theta - t\sin\theta, s\sin\theta + t\cos\theta) dt.
\]
This is the \textbf{Radon transform\index{Radon transform}} $\mathcal{R}f(\theta, s)$, \cite{foschiatti2025tomography}.

\textbf{Inversion} (filtered back-projection):
\[
f(x,y) = \int_0^\pi \left[ p_\theta * h \right](x\cos\theta + y\sin\theta) d\theta,
\]
where $h$ is the \textbf{Ram-Lak filter} (ramp filter in Fourier domain): $\hat{h}(\omega) = |\omega|$.

%==============================================================================
\section{Experimental Setup: Visible-Light Tomography}
\label{sec:tom:setup}

\textbf{Hardware}:
\begin{itemize}
    \item Arduino-controlled stepper motor (rotation stage)
    \item Projector (visible light source) + Fresnel lens (collimation)
    \item Camera (detector) + white screen
    \item Transparent cellulose acetate sheets with printed text (phantom scroll)
\end{itemize}

\textbf{Acquisition}:
\begin{lstlisting}[style=python,caption=Arduino stepper control + camera sync]
# Arduino sketch (simplified)
#include <Stepper.h>
const int stepsPerRevolution = 200;
Stepper myStepper(stepsPerRevolution, 8, 9, 10, 11);
int nProjections = 800;

void setup() {
    myStepper.setSpeed(10);  // rpm
    Serial.begin(9600);
}

void loop() {
    if (Serial.available()) {
        char cmd = Serial.read();
        if (cmd == 's') {  // start acquisition
            for (int i = 0; i < nProjections; i++) {
                myStepper.step(stepsPerRevolution / nProjections);
                delay(500);  // settle
                Serial.println("CAPTURE");  // trigger camera
                delay(500);
            }
        }
    }
}
\end{lstlisting}

%==============================================================================
\section{Reconstruction Pipeline}
\label{sec:tom:pipeline}

\subsection{1. Preprocessing}
\begin{itemize}
    \item Dark/flat field correction
    \item Rotation axis determination (cross-correlation of $0^\circ$ and $180^\circ$ projections)
    \item Cropping to region of interest
\end{itemize}

\subsection{2. Filtered Back-Projection (FBP)}
\begin{lstlisting}[style=python,caption=FBP reconstruction using iradon]
import numpy as np
from skimage.transform import iradon

# projections: (n_angles, n_pixels) array
# angles: in degrees
reconstruction = iradon(projections, theta=angles, 
                        filter='ramp',  # Ram-Lak
                        interpolation='linear',
                        circle=False)

# Post-process: circular mask to remove edge artifacts
from skimage.draw import disk
mask = np.zeros_like(reconstruction)
rr, cc = disk((n/2, n/2), n/2 - 5)
mask[rr, cc] = 1
reconstruction *= mask
\end{lstlisting}

\subsection{3. Segmentation: Finding Layers}
Each slice is a cross-section through the scroll. The layers appear as concentric curved lines.

\begin{lstlisting}[style=python,caption=Spiral segmentation via genetic algorithm]
import numpy as np
from scipy.optimize import differential_evolution

def spiral_model(theta, params):
    # params = [center_x, center_y, a, b]
    cx, cy, a, b = params
    r = a + b * theta
    return cx + r * np.cos(theta), cy + r * np.sin(theta)

def segmentation_loss(params, slice_img):
    # Generate spiral mask from params
    # Compare to image gradients
    # Return negative correlation
    pass

# Optimize per slice
result = differential_evolution(segmentation_loss, bounds,
                                args=(slice_img,), maxiter=100)
best_spiral = result.x
\end{lstlisting}

\subsection{4. Texture Mapping and Flattening}
For each layer, extract intensity values along the spiral, map to a rectangular image.

\begin{lstlisting}[style=python,caption=Layer flattening]
def flatten_layer(volume, spiral_params, layer_thickness=3):
    """Extract and flatten one layer from 3D volume."""
    h, w, d = volume.shape
    flattened = np.zeros((d, n_points))
    
    for z in range(d):
        slice_img = volume[:, :, z]
        cx, cy, a, b = spiral_params[z]
        
        # Sample along spiral
        theta = np.linspace(0, 4*np.pi, n_points)
        x = cx + (a + b*theta) * np.cos(theta)
        y = cy + (a + b*theta) * np.sin(theta)
        
        # Interpolate
        from scipy.ndimage import map_coordinates
        flattened[z, :] = map_coordinates(slice_img, [y, x], order=1)
    
    return flattened
\end{lstlisting}

\subsection{5. Maximum Intensity Projection (MIP)}
Combine multiple layers into one readable image:
\[
\mathrm{MIP}(u,v) = \max_{z} I_z(u,v).
\]
Invert if text is dark-on-light.

%==============================================================================
\section{Results: Reconstructing the Herculaneum Scrolls} \cite{foschiatti2025tomography}
\label{sec:tom:results}

The article demonstrates reconstruction of:
\begin{itemize}
    \item \textbf{Greek}: Homer, \textit{Hymni Homerici XV} (Hercules hymn)
    \item \textbf{Latin}: Lucretius, \textit{De Rerum Natura} (Book 1, lines 72--74)
\end{itemize}

%\begin{figure}[ht]
%\centering
%\includegraphics[width=0.8\textwidth]{figures/greek_reconstruction.png}
%\caption{Greek scroll reconstruction: original (left), MIP reconstruction (right).}
%\end{figure}

%==============================================================================
\section{Code Repository: Archeolab}
\label{sec:tom:archeolab}

Complete pipeline (MATLAB + Python) at:
\begin{center}
\url{https://gitlab.com/csc1/archeolab/}
\end{center}

Includes:
\begin{itemize}
    \item \texttt{guiAxisFinder.m} --- rotation axis detection
    \item \texttt{guiSegment.m} --- manual + genetic algorithm segmentation
    \item \texttt{FBP\_reconstruction.m} --- filtered back-projection
    \item \texttt{MIP\_flattening.m} --- texture mapping + flattening
    \item Jupyter notebooks for Python equivalents
    \item Pre-acquired datasets for testing without hardware
\end{itemize}

%==============================================================================
\section{Bridge to Chapter 14}
\label{sec:tom:bridge}

Tomography reconstructs \emph{hidden structure} from projections. Chapter 14 shows how \emph{FAIR principles} reconstruct \emph{usable knowledge} from legacy mathematical software. Both are inverse problems: from fragmented observations to coherent understanding.

%==============================================================================
\section{Further Reading}
\label{sec:tom:refs}

\begin{itemize}
    \item Foschiatti, Kittenberger, Scherzer (2025). \textit{Deciphering Scrolls with Tomography}. Notices, 72(10).
    \item Seales et al. (2016). \textit{From Damage to Discovery via Virtual Unwrapping}. Science Advances.
    \item Natterer (2001). \textit{The Mathematics of Computerized Tomography}. SIAM.
    \item Feeman (2015). \textit{The Mathematics of Medical Imaging}. Springer.
    \item Vesuvius Challenge: \url{https://vesuviuschallenge.org}
\end{itemize}

%==============================================================================
\section{Exercises}
\label{sec:tom:exercises}

\begin{exercise}
Implement the Radon transform and its adjoint (back-projection) from scratch in NumPy. Verify $\mathcal{R}^*\mathcal{R}$ is a convolution with $1/|x|$, \cite{foschiatti2025tomography}.
\end{exercise}

\begin{exercise}
Build the visible-light tomography setup (projector + camera + Arduino stage). Acquire 180 projections of a printed spiral phantom and reconstruct with FBP, \cite{foschiatti2025tomography}.
\end{exercise}

\begin{exercise}
The Ram-Lak filter amplifies high-frequency noise. Implement a \emph{windowed} ramp filter (Hann, Hamming, Shepp-Logan) and compare reconstruction quality vs. noise amplification.
\end{exercise}

\begin{exercise}[Research]
The Herculaneum scrolls have carbon-based ink (low X-ray contrast). Develop a phase-contrast CT simulation and reconstruction pipeline. How much does phase retrieval improve ink visibility? \cite{foschiatti2025tomography}
\end{exercise}